\section{SDM continua y explicación de los componentes del Transformer}

La SDM clásica usa binary address y hard radius; el Transformer usa continuous vector y softmax. A continuación, la clase vuelve continua la SDM y plantea una pregunta más abierta: si un Transformer entrenado cae realmente en un «intervalo de parámetros de SDM útil», y si la FFN, la LayerNorm y la memory capacity pueden integrarse en este mismo marco.

\subsection{Del binary radius al kernel continuo}

La diapositiva de Continuous SDM explica que no es necesario cuantizar la entrada en bits. Para vectores reales normalizados, basta con usar directamente un dot-product kernel exponencial para definir una lectura continua:
$$
\operatorname{CSDM}(q)=
\sum_\mu\operatorname{softmax}_\mu(\beta p_\mu^{\top}q)v_\mu.
$$
Esto tiene la misma forma que la attention y también puede implementarse con un perceptrón multicapa o con una associative-memory network. La versión continua conserva el principio de que «las direcciones cercanas contribuyen más», pero convierte el radio duro en pesos suaves.

\lecturefigure{slide-31-continuous-sdm.jpg}{Continuous SDM: del binary Hamming space a vectores reales}{Diapositiva V031 recuperada de la grabación oficial; 00:21:32}

\begin{knowledgebox}{Qué es la effective beta}
$\beta$ determina la retrieval sharpness: una $\beta$ pequeña hace que muchas memories participen promediadas; una $\beta$ grande hace que domine la memory más similar. En el Transformer, la scale explícita, las key/query norms y la LayerNorm modifican conjuntamente la temperatura efectiva, por lo que para medir el learned coefficient hay que observar los logits reales y no solo la constante de la fórmula.
\end{knowledgebox}

\subsection{¿Se parece la Attention entrenada a una SDM «buena»?}

La diapositiva de la pregunta propone una prueba inversa: si la attention implementa una SDM, ¿corresponden los parámetros entrenados a un radio/temperature con buena capacidad y buena robustez al ruido? La diapositiva de Learned beta muestra la distribución del coefficient en distintas heads, lo que indica que el modelo no tiene una temperatura uniforme. Cada head puede encargarse de tareas distintas, como copiar, sintaxis local o recuperación de largo alcance, por lo que la retrieval sharpness óptima también puede diferir.

\lecturefigure{slide-32-transformer-sdm-question.jpg}{Pregunta abierta: ¿corresponde un Transformer entrenado a parámetros de SDM óptimos?}{Diapositiva V032 recuperada de la grabación oficial; 00:22:18}

\lecturefigure{slide-33-learned-beta.jpg}{La effective beta aprendida por distintas attention heads}{Diapositiva V033 recuperada de la grabación oficial; 00:23:44}

\begin{warningbox}{No tomar la entropy como indicador monótono de calidad}
Una entropy alta puede indicar que hace falta integrar varias memorias, pero también que la recuperación no está enfocada; una entropy baja puede ser una recuperación precisa, pero también un exceso de confianza en un error. La perspectiva de la SDM ofrece una explicación en términos de capacidad/ruido, pero el juicio debe combinarse con la función de la head y la pérdida de la tarea.
\end{warningbox}

\subsection{FFN, LayerNorm y el memory framework}

La diapositiva de Transformer components integra la broader architecture en el marco de la SDM: la FFN puede verse como un conjunto de key--value memories de largo plazo; la LayerNorm ayuda a estabilizar la norma de los vectores y el coefficient; la attention es una pattern-based retrieval explícita. Por último, la diapositiva de summary presenta las conclusiones junto a las preguntas futuras y subraya que estas correspondencias pueden generar hipótesis de investigación, pero no constituyen una teoría unificada ya demostrada.

\lecturefigure{slide-34-transformer-components-converge.jpg}{Cómo convergen los componentes del Transformer en el SDM framework}{Diapositiva V034 recuperada de la grabación oficial; 00:25:06}

\lecturefigure{slide-35-sdm-summary.jpg}{Resumen de la parte teórica de la primera exposición y preguntas abiertas}{Diapositiva V035 recuperada de la grabación oficial; 00:26:20}

\begin{longtable}{@{}L{0.20\textwidth}L{0.31\textwidth}L{0.37\textwidth}@{}}
\toprule
Componente del Transformer & Perspectiva de SDM/memoria & Limitación que debe conservarse \\
\midrule
Attention & Pattern-based retrieval precisa en el contexto actual & Restringida por la context window; almacena todas las keys/values \\
FFN & Key--value associations de largo plazo superpuestas en los parámetros & Recuperación más ruidosa y más difícil de rastrear elemento por elemento; no equivale a almacenar el conjunto de datos completo \\
LayerNorm & Estabiliza la norm de los vectores para que la correspondencia distancia/producto punto sea más controlable & Las norms reales y las learned scales siguen modificando la temperature \\
Multi-head & Varios subespacios de representación distintos y distintas retrieval sharpness & No puede equipararse directamente con las microzones de las regiones cerebrales \\
\bottomrule
\end{longtable}

\teachervoice{En la sesión de preguntas, el ponente distingue entre la SDM neuron-centric y la pattern-centric: la attention conserva todas las pattern locations, por lo que la similitud es precisa pero existe una context window; la memoria distribuida de largo plazo solo conserva una noisy superposition, con mayor capacidad pero con dificultad para triangular con precisión el pattern original. Este tradeoff se acerca más a la verdadera enseñanza de ingeniería de la clase que las métricas de monitoreo en producción que inventaba la versión anterior.}

\subsection{Resumen de la sección}

La SDM continua convierte el radio duro en un kernel exponencial y comparte su forma con la attention. Las explicaciones de la learned beta, la LayerNorm y la FFN extienden el marco al Transformer completo, pero cada correspondencia debe tratarse como una hipótesis comprobable. La precisión en contextos cortos de la pattern memory frente a la alta capacidad ruidosa de la distributed memory es el compromiso central que plantea la clase.
